% p-glosario.tex
% !TEX root = z-main.tex

\newglossaryentry{3d}{
    name={3D},
    description={Término que se refiere a representaciones tridimensionales, es decir, con ancho, alto y profundidad}
}

\newglossaryentry{RA}{
    name={RA},
    description={Acrónimo de Realidad Aumentada, tecnología que superpone elementos digitales en el entorno físico del usuario}
}

\newglossaryentry{RV}{
    name={RV},
    description={Acrónimo de Realidad Virtual, tecnología que sumerge al usuario en un entorno digital completamente generado por computadora}
}

\newglossaryentry{ARCore}{
    name={ARCore},
    description={Plataforma de desarrollo de realidad aumentada creada por Google, que permite construir aplicaciones que integran elementos digitales en el mundo real}
}

\newglossaryentry{Wireframe}{
    name={Wireframe},
    description={Representación visual de un modelo tridimensional utilizando líneas y aristas para mostrar su estructura básica sin superficies sólidas}
}

\newglossaryentry{pipeline}{
    name={pipeline},
    description={Conjunto de procesos y etapas secuenciales que se siguen para lograr un objetivo específico, como la creación de modelos 3D o el desarrollo de software}
}


\newglossaryentry{PBR}{
    name={PBR},
    description={Acrónimo de Physically Based Rendering, técnica de renderizado que simula el comportamiento físico de la luz para lograr representaciones más realistas de materiales y superficies}
}